\documentclass[10pt]{beamer}
\usetheme{PaloAlto}
\usecolortheme{seahorse}
\setbeamertemplate{navigation symbols}{}
\setbeamertemplate{caption}[numbered]
%general package
\usepackage[utf8]{inputenc}
\usepackage{array}
\usepackage[english]{babel}
\usepackage{geometry}
\usepackage{tcolorbox}
\usepackage[export]{adjustbox}
\usepackage{graphicx}
\usepackage{xcolor}
\graphicspath{{../img}}
\hypersetup{
    colorlinks=true, 
    linkcolor=red,    % Internal links
    urlcolor=blue     % External links
}
\usepackage{wrapfig}
\usepackage{luacode}
\usepackage{comment}

%math package
\usepackage{amsmath}
\usepackage{amsfonts}
%\usepackage{amssymb}
%\usepackage{amsthm}
%\usepackage{slashed}
%\usepackage{tikz-cd}

%font package
%\usepackage{mathrsfs}
%\usepackage{bm}

%misc. package
\usepackage{enumitem}
\usepackage{animate}

\DeclareMathOperator{\xd}{\,d\!}
\DeclareMathOperator{\curl}{curl}
\DeclareMathOperator{\dive}{div}
\newcommand{\e}{{\rm e}}
\newcommand{\norm}[1]{\lVert#1\rVert}
\newcommand{\R}{\mathbb R}
\newcommand{\vF}{\mathbf F}
\newcommand{\vv}{\mathbf v}
\newcommand{\inpr}[1]{\left\langle#1\right\rangle}

\author[B.H.]{{\Large A Collection of Daily Exercises}\\\vspace{.5em} Author: Benjamin Huang}
\date{}
\title[Skill Builder]{Math Skill Builder}

%\institute[]{\vskip -2em\includegraphics[width = 0.65\textwidth]{}}
%\logo{\includegraphics[width=.12\textwidth]{}}

\begin{document}
\begin{luacode*}
function generateNonzero(min, max)
  local output
  repeat
    output = math.random(min, max)
  until output ~= 0
  return output
end
\end{luacode*}
\frame{\titlepage}
\begin{frame}
\frametitle{Trigonometric Function Properties Review}
\[
A\sin(bx) + c
\]
\pause
\begin{itemize}
\item
$\displaystyle \frac{2\pi}{b}$ is called the \textbf{period}.\pause
\item
$y = c$ is called the \textbf{midline}.\pause
\item
$\lvert A\rvert$ is called the \textbf{amplitude}.\pause
\item
The graph of this function is called a \textbf{sinusoidal curve}.\pause
\item 
Same terminology applies to the cosine functions either.\pause
\end{itemize}
Miscellaneous concepts:\pause
\begin{itemize}
\item
$b$ is called the \textbf{angular speed} or \textbf{angular frequency}.\pause
\item
$\displaystyle \frac{b}{2\pi}$ is called the \textbf{frequency}.
\end{itemize}
\end{frame}
\end{document}